\documentclass{article}
\usepackage{amsmath, amssymb}
\usepackage{amsthm}

\author{Michael Naumov}
\date{}
\setlength{\parindent}{0pt}
\setlength{\parskip}{0.5em}


\title{1.2.4-21}

\begin{document}
\maketitle

[1.2.1-5](<../1.2.1. Mathematical induction/1.2.1-5.md>)

\( \forall n > 1: n = p_1p_2 \cdot\ \cdots\ \cdot p_k, p_1 \leq p_2 \leq \cdots \leq p_k \)

Assume that not all \( n \) have such an unique expansion.

Let \( n \) be the smallest n with non-unique expansion.

\( n = p_1p_2 \cdot\ \cdots\ \cdot p_k = q_1q_2 \cdot\ \cdots\ \cdot q_l, p_1 \leq p_2 \leq \cdots p_k, q_1 \leq q_2 \leq \cdots \leq q_l \)

Where \( (p_1,p_2,\ldots,p_k) \neq (q_1,q_2,\ldots,q_l) \)

\section*{1}

If \( p_1 = q_1 \)

\( m = \frac{n}{p_1} = \frac{n}{q_1} < n \)

\( m = p_2 \cdot\ \cdots\ \cdot p_k = q_2 \cdot\ \cdots\ \cdot q_l \)

\( m < n \implies m \) has an unique expansion \( \implies k = l, p_j = q_j, 2 \leq j \leq k \)

And also \( p_j = q_j \) for \( j = 1 \)

\( \implies (p_1,p_2,\ldots,p_k) = (q_1,q_2,\ldots,q_l) \)

Contradiction

\section*{2}

If \( p_1 \neq q_1 \).

We can assume \( p_1 < q_1 \)

Otherwise we just change notations \( (p_1, p_2, \ldots, p_k) \leftrightarrow (q_1, q_2, \ldots, q_l) \)

\( \forall j, 1 \leq j \leq l: q_j > p_1 \)

Let \( d = \gcd(q_j, p_1) \)

\( d \backslash  q_j \)

\( q_j \) - prime \( \implies d = 1 \lor d = q_j \)

\( q_j > p_1 \implies \lnot(q_j \backslash p_1) \implies d = 1 \)

\( \implies q_j \perp p_1 \)

\( q_1q_2 \cdot\ \cdots\ \cdot q_l = n \equiv 0 = q_1 \cdot 0\ (\mathrm{modulo}\ p_1) \)

Law B \( \implies q_2 \cdot\ \cdots\ \cdot q_l \equiv 0 = q_2 \cdot 0\ (\mathrm{modulo}\ p_1) \)

Law B \( \implies q_3 \cdot\ \cdots\ \cdot q_l \equiv 0 = q_3 \cdot 0\ (\mathrm{modulo}\ p_1) \)

\( \vdots \)

\( q_l = q_l \cdot 1 \equiv 0 = q_l \cdot 0\ (\mathrm{modulo}\ p_1) \)

Law B \( \implies 1 \equiv 0\ (\mathrm{modulo}\ p_1) \implies p_1 \backslash 1 \) - Contradiction


\end{document}
